\subsection{Convex Hull}

The convex hull of a point set $S \subset \mathbb{R}^2$ is the smallest convex
set containing $S$:
\[
  \text{conv}(S) = \left\{ \sum_i \lambda_i p_i : p_i \in S,\; \lambda_i \geq 0,\;
  \sum_i \lambda_i = 1 \right\}.
\]
For a polygon $D$, $\text{conv}(D)$ is determined by its vertex set.

\subsection{Convex Hull Ratio}

\begin{equation}
  \text{CH}(D) = \frac{A(D)}{A(\text{conv}(D))} \in (0, 1].
  \label{eq:ch}
\end{equation}

\begin{proposition}
  CH$(D) = 1$ if and only if $D$ is convex.
\end{proposition}

\begin{proof}
  If $D$ is convex, then $D = \text{conv}(D)$, so $A(D) = A(\text{conv}(D))$
  and CH $= 1$. Conversely, if $D$ is non-convex, there exists a point $p$
  outside $D$ but inside $\text{conv}(D)$, so $A(\text{conv}(D)) > A(D)$
  and CH $< 1$. \qed
\end{proof}

\subsection{Canonical Values}

\begin{table}[ht]
\centering
\caption{Convex Hull Ratio for canonical shapes.}
\label{tab:ch_canonical}
\begin{tabular}{lc}
\toprule
Shape & CH \\
\midrule
Any convex polygon & $1.00$ \\
Regular star polygon ($n=5$) & $\approx 0.38$ \\
Cross shape ($L = 10$, $w = 2$) & $0.80$ \\
Corridor district (width $= 0.5$, length $= 10$) & $\approx 0.50$ \\
L-shaped district & $\approx 0.75$ \\
\bottomrule
\end{tabular}
\end{table}

\subsection{Relationship to PP and Reock}

Unlike PP and Reock, CH has no simple algebraic relationship to other metrics.
CH measures specifically the degree of non-convexity, which is a distinct
geometric property from:
\begin{itemize}
  \item Perimeter raggedness (PP, S): a smooth non-convex boundary can have
    low CH but moderate PP.
  \item Global circularity (Reock): a district can be elongated (low Reock)
    but convex (CH $= 1$), or compact globally (high Reock) but have thin arms
    (low CH).
\end{itemize}
This independence from PP and Reock makes CH a non-redundant component of
the multi-metric compactness profile.
